\documentclass[11pt]{article}
\usepackage[a4paper,margin=1in]{geometry}
\usepackage{amsmath,amssymb,amsthm}
\usepackage{booktabs}
\usepackage{enumitem}
\usepackage{hyperref}
\usepackage{mathtools}
\usepackage{bm}

\title{Algorithmic Improvement Proposals for Vamana ANN\\
\large A Rigorous, Implementation-Oriented Exploration}
\author{Prepared for the GraphANN repository}
\date{\today}

\begin{document}
\maketitle

\begin{abstract}
This document proposes five high-impact algorithmic improvements for a baseline Vamana-style approximate nearest neighbor (ANN) implementation. For each proposal, we provide: (i) concrete changes to the current algorithm and code structure, (ii) expected empirical benefits in recall/latency/build-time tradeoffs, and (iii) rigorous mathematical framing to justify why the change should improve graph navigability under fixed memory budgets. The proposals are prioritized for practical adoption in an educational C++ implementation that already supports greedy search, $\alpha$-RNG pruning, bounded degree, and parallel insertion.
\end{abstract}

\section{Baseline Summary and Optimization Goal}

Let $P=\{x_i\}_{i=1}^n \subset \mathbb{R}^d$ denote the dataset and let $G=(V,E)$ be a directed graph over points with bounded out-degree. Query-time ANN search uses best-first greedy graph traversal with candidate budget $L$, returning top-$K$ points by distance.

The practical optimization objective is a constrained multi-objective tradeoff:
\begin{equation}
\min_{G,\,\text{search policy}} \Big(\mathbb{E}[\text{latency}],\; \mathbb{E}[\text{distance computations}],\; \text{build time}\Big)
\quad\text{s.t.}\quad \text{Recall@}K \ge \rho,\; \deg^+(v) \le R.
\end{equation}

Vamana-style design specifically tries to construct a sparse graph that remains highly navigable under greedy traversal. The five improvements below target the same constrained objective through better graph quality, better candidate generation, and lower per-query compute.

\section{Improvement 1: Two-Pass Vamana with Medoid Start and $\alpha$ Scheduling}

\subsection{What to change}
\begin{enumerate}[leftmargin=*]
    \item Replace random entry node with a medoid-like center:
    \begin{equation}
    \mu = \frac{1}{n}\sum_{i=1}^n x_i,\qquad s = \arg\min_i \lVert x_i - \mu \rVert_2^2.
    \end{equation}
    \item Initialize graph with random neighbors (small over-connected random seed) instead of fully empty adjacency.
    \item Perform \textbf{two full insertion passes}:
    \begin{itemize}[leftmargin=*]
        \item Pass 1: strict pruning with $\alpha_1 = 1.0$.
        \item Pass 2: relaxed pruning with $\alpha_2 > 1$ (typical range: $1.1$--$1.4$).
    \end{itemize}
    \item Keep degree controls (target $R$, overflow threshold $\gamma R$), optionally with pass-specific $L$ values.
\end{enumerate}

\subsection{Why it should be better}
\begin{itemize}[leftmargin=*]
    \item A random start node increases expected initial search distance and can bias early graph shaping.
    \item A medoid-like entry improves average routing quality from the first hop.
    \item Single-pass construction is path-dependent; two-pass refinement reduces local construction artifacts.
    \item $\alpha_1=1$ preserves local geometric precision; $\alpha_2>1$ reintroduces strategic long-range links that reduce diameter/hop count.
\end{itemize}

\subsection{Mathematical perspective}
Define expected hop count to an acceptable neighborhood of true nearest neighbors as $H(q;G,s)$. Under fixed degree budget, we seek:
\begin{equation}
\min_G\; \mathbb{E}_q[H(q;G,s)]\quad\text{s.t.}\quad \deg^+(v)\le R.
\end{equation}

Interpret two-pass Vamana as continuation in pruning aggressiveness:
\begin{equation}
G^{(1)} = \mathcal{P}_{\alpha=1}(G_0),\qquad G^{(2)} = \mathcal{P}_{\alpha=\alpha_2}(G^{(1)}),\quad \alpha_2>1.
\end{equation}

Pass 1 sharpens local structure; pass 2 broadens navigability. Empirically this often decreases graph diameter and high-percentile query hops while keeping degree bounded.

\subsection{Cost impact}
Build time typically increases by roughly $1.7\times$--$2.0\times$, but query-time $L$ for a target recall can often be reduced.

\section{Improvement 2: Candidate Set for Prune from Visited Set (Not Only Final Frontier)}

\subsection{What to change}
During insertion of point $p$, greedy search should expose two sets:
\begin{itemize}[leftmargin=*]
    \item Final bounded frontier $S_L(p)$ (current behavior).
    \item Full visited/expanded set $V_p$ (new required tracking).
\end{itemize}

Then build prune candidates as:
\begin{equation}
C_p = V_p \cup N_{\text{out}}(p)
\end{equation}
instead of using only final frontier elements. Run robust prune on $C_p$ to select up to $R$ neighbors.

\subsection{Why it should be better}
\begin{itemize}[leftmargin=*]
    \item The final frontier is a compressed, lossy summary of traversal.
    \item Many useful bridge nodes are visited but not retained in top-$L$ frontier.
    \item Richer candidate support improves odds of selecting edges that reduce local minima traps and improve global routing.
\end{itemize}

\subsection{Mathematical perspective}
Let $\mathcal{F}(A)$ denote the utility of selected neighbors under proximity+diversity constraints induced by robust pruning. Current optimization is over a restricted domain:
\begin{equation}
\max_{A\subseteq S_L(p),\,|A|\le R} \mathcal{F}(A).
\end{equation}

Improved optimization uses the larger domain:
\begin{equation}
\max_{A\subseteq (V_p\cup N_{\text{out}}(p)),\,|A|\le R} \mathcal{F}(A).
\end{equation}

Since $S_L(p)\subseteq V_p\cup N_{\text{out}}(p)$, the optimum value cannot decrease.

\subsection{Cost impact}
Prune stage costs increase because candidate list is larger. However, better graph quality can reduce required query beam width for the same recall target.

\section{Improvement 3: Locally Adaptive $\alpha_i$ and Degree Budget $R_i$}

\subsection{What to change}
\begin{enumerate}[leftmargin=*]
    \item Estimate local intrinsic dimensionality (LID) per point $i$ from local distance statistics:
    \begin{equation}
    \widehat{m}_i = \left(-\frac{1}{k-1}\sum_{j=1}^{k-1}\log\frac{r_j(i)}{r_k(i)}\right)^{-1},
    \end{equation}
    where $r_j(i)$ is the $j$-th nearest local sample distance.
    \item Set node-specific pruning aggressiveness and degree:
    \begin{equation}
    \alpha_i = \mathrm{clip}\left(1+\frac{c_\alpha}{\sqrt{\widehat{m}_i}},\;\alpha_{\min},\alpha_{\max}\right),
    \end{equation}
    \begin{equation}
    R_i = \mathrm{clip}\left(R_0 + c_R\widehat{m}_i,\;R_{\min},R_{\max}\right).
    \end{equation}
    \item Replace global prune parameters with per-node values in robust prune and overflow logic ($\gamma R_i$ threshold per node).
    \item Keep total memory bounded by enforcing
    \begin{equation}
    \sum_{i=1}^n R_i \le nR_{\text{budget}}.
    \end{equation}
\end{enumerate}

\subsection{Why it should be better}
\begin{itemize}[leftmargin=*]
    \item Real datasets are locally heterogeneous.
    \item Uniform $\alpha$ and $R$ over-prune hard regions and over-allocate easy regions.
    \item Adaptive budgets allocate edges where marginal recall gain per edge is larger.
\end{itemize}

\subsection{Mathematical perspective}
Global policy solves:
\begin{equation}
\min_G \sum_i \mathcal{L}_i(G;\alpha,R).
\end{equation}
Adaptive policy solves:
\begin{equation}
\min_G \sum_i \mathcal{L}_i(G;\alpha_i,R_i)
\quad\text{s.t.}\quad
\sum_i R_i \le nR_{\text{budget}}.
\end{equation}

This is a constrained resource allocation problem. LID is a proxy for local geometric complexity and expected branching needs.

\subsection{Cost impact}
Small overhead for LID estimation and metadata storage; potentially strong recall-latency gains at fixed average memory.

\section{Improvement 4: Multi-Entry Search and Query-Adaptive Beam Width}

\subsection{What to change}
\begin{enumerate}[leftmargin=*]
    \item Replace single start node with $M$ entry points (medoid + diversified pivots/centroids).
    \item Seed search frontier with all entries (shared best-first traversal).
    \item Use adaptive beam width:
    \begin{itemize}[leftmargin=*]
        \item start with small $L_0$,
        \item increase only when confidence is low or progress stalls.
    \end{itemize}
    \item Add confidence-based stopping criterion using candidate lower bounds and current $k$-th best distance.
\end{enumerate}

\subsection{Why it should be better}
\begin{itemize}[leftmargin=*]
    \item Tail queries far from the single entry basin are brittle.
    \item Multiple entries reduce local minima risk.
    \item Adaptive $L$ avoids over-compute on easy queries and spends work only on hard ones.
\end{itemize}

\subsection{Mathematical perspective}
If one entry fails with probability $p_f$ and starts are weakly correlated, multi-entry failure behaves as:
\begin{equation}
 p_f^{(M)} \approx p_f^M.
\end{equation}

For stopping, let $d_k$ be current exact $k$-th best and $\ell_{\text{next}}$ a lower bound on unexpanded candidates. Stop when:
\begin{equation}
\ell_{\text{next}} \ge d_k(1-\varepsilon),
\end{equation}
which bounds expected gain from additional expansions.

\subsection{Cost impact}
Low metadata overhead for extra entries; often large reduction in high-percentile latency at constant recall.

\section{Improvement 5: Distance Cascade (Approximate Screening + Exact Rerank)}

\subsection{What to change}
\begin{enumerate}[leftmargin=*]
    \item Store compact approximate representation per point (e.g., scalar-quantized or PQ-like code).
    \item In greedy expansion and prune pairwise checks, compute cheap approximate distance $\tilde d$ first.
    \item Compute exact distance $d$ only for top-$B$ approximate candidates or ambiguous margin cases.
    \item Keep final top-$K$ reranking exact to preserve output quality.
\end{enumerate}

\subsection{Why it should be better}
Distance computation dominates CPU cost in both build and query in graph ANN systems. Most candidates are clearly inferior and can be rejected with cheap approximate scores.

\subsection{Mathematical perspective}
Assume bounded approximation error:
\begin{equation}
|d(x,y)-\tilde d(x,y)|\le \delta.
\end{equation}
If candidates $a,b$ satisfy
\begin{equation}
\tilde d(a)+\delta < \tilde d(b)-\delta,
\end{equation}
then exact ordering is guaranteed ($d(a)<d(b)$), so one exact evaluation can be skipped safely. This yields a reduction factor:
\begin{equation}
\mathbb{E}[C_{\text{exact}}] \approx \rho\,\mathbb{E}[C_{\text{total}}],\qquad \rho\ll 1,
\end{equation}
while maintaining quality through exact final reranking.

\subsection{Cost impact}
Additional memory and calibration effort, but often the largest raw speedup in CPU-bound regimes.

\section{Implementation Priority and Experimental Plan}

\subsection{Recommended order}
\begin{enumerate}[leftmargin=*]
    \item Two-pass build + medoid start.
    \item Visited-set candidate generation for robust prune.
    \item Multi-entry + adaptive beam search.
    \item Adaptive per-node $\alpha_i, R_i$.
    \item Distance cascade with exact rerank.
\end{enumerate}

\subsection{Ablation protocol}
For each stage, benchmark on fixed datasets and report:
\begin{itemize}[leftmargin=*]
    \item Recall@10, Recall@100.
    \item Mean and P99 latency.
    \item Mean distance computations/query.
    \item Build time and memory footprint.
    \item Degree distribution and graph diameter proxy (or average hop count to top-$K$ basin).
\end{itemize}

Use one-change-at-a-time ablations and then cumulative composition to detect interactions.

\section{Expected Outcomes}

Under fixed average degree budget, these five modifications are expected to:
\begin{itemize}[leftmargin=*]
    \item improve graph navigability (especially tail queries),
    \item reduce required beam width for target recall,
    \item lower high-percentile latency,
    \item preserve or improve recall under comparable memory,
    \item provide a stronger research platform for further ANN experiments.
\end{itemize}

\section*{Conclusion}
The proposed roadmap preserves the spirit of Vamana while addressing core limitations of the current baseline: construction path dependence, candidate impoverishment, parameter rigidity, entry-point brittleness, and expensive exact distance overuse. The combination of graph-quality upgrades (Improvements 1--3) and search-cost controls (Improvements 4--5) gives a principled path to better recall-latency-memory tradeoffs in practice.

\end{document}
